\documentclass[11pt]{article}
\usepackage[T1]{fontenc}
\usepackage{textcomp}
\usepackage[margin=0.75in]{geometry}
\usepackage{amsmath,amssymb,amsthm}
\usepackage{array,booktabs}
\usepackage{hyperref}
\setlength{\parindent}{0pt}
\newtheorem{theorem}{Theorem}
\theoremstyle{definition}
\newtheorem{definition}{Definition}
\title{Flow Proof Artifact: prop-10-circles-touch-at-most-one-point}
\author{Generated by \texttt{flow doc proof}}
\date{Flow Proof Book}
\begin{document}
\maketitle
One circle cannot touch another at more than one point.\\[0.5em]
\textbf{Source.} euclid — Elements, Book III, Proposition 10\\[0.5em]
\bigskip
\noindent{\large \textbf{Derived fact 1.} \textit{One circle cannot touch another at more than one point} \hfill the Euclidean plane \cdot Euclid Book III \cdot Proposition 10: one circle cannot touch another at more than one point}\par
\smallskip\noindent\textit{Coordinate: the Euclidean plane \cdot Euclid Book III \cdot Proposition 10: one circle cannot touch another at more than one point}\par
\smallskip\noindent\textit{Source: euclid — Elements, Book III, Proposition 10}\par
\smallskip\noindent\textit{Built on: proposition 6: the line joining centers passes through the point of contact, for Euclid Book III on the Euclidean plane}\par
\medskip
\noindent\textbf{Goal.} One circle cannot touch another at more than one point\par
\noindent$\displaystyle \text{touching circles have unique contact}$\par
\medskip
\noindent
\begin{tabular}{@{} >{\raggedright\arraybackslash}p{0.06\textwidth} >{\raggedright\arraybackslash}p{0.47\textwidth} >{\raggedleft\arraybackslash}p{0.41\textwidth} @{}}
\textbf{\#} & \textbf{Proof} & \textbf{Mathematics} \\
\midrule
\textbf{1.} & We prove that proposition 10: one circle cannot touch another at more than one point for Euclid Book III the Euclidean plane. & \\[0.45em]
\textbf{2.} & Let C1 = first\_circle. & \\[0.45em]
\textbf{3.} & Let C2 = second\_circle. & \\[0.45em]
\textbf{4.} & Let T = point\_of\_contact. & \\[0.45em]
\textbf{5.} & From step 2, step 3, and step 4, we can deduce that at most one contact point. & $\displaystyle \text{at most one contact point}$ \\[0.45em]
\textbf{6.} & We invoke the derived fact governing Euclid Book III the Euclidean plane: proposition 6: the line joining centers passes through the point of contact, for Euclid Book III on the Euclidean plane. & \\[0.45em]
\textbf{7.} & From step 6, this implies touching circles have unique contact. Hence proven. & $\displaystyle \text{touching circles have unique contact}$ \\[0.45em]
\bottomrule
\end{tabular}
\label{thm:Geometry-Euclid Book III-proposition_10_one_circle_cannot_touch_another_at_more_than_one_point}
\par\medskip\hrule
\end{document}
